\section{Binding abstract resources to apparatus variables}
\label{sec:physical-binding}

The theorem layer above is intentionally agnostic about what one unit of $e_k$ means. Resource closure requires an explicit map from apparatus variables to the per-unit target and nuisance scores, the covariance model, and the cost vector. RQIR therefore separates the \emph{allocation theorem} from the \emph{apparatus binding}.

\subsection{Shots and repeated trials}

For conditionally independent repeated trials at setting $k$, the simplest binding is
\begin{equation}
 e_k=N_k,
\end{equation}
where $N_k$ is the number of accepted trials and $(\bm s_k,J_k)$ are per-trial whitened scores. The continuous design problem provides an approximate lower bound when $N_k$ must be integer. Dead time, preparation failure, postselection, or state-reset overhead must either modify the accepted-trial rate or appear in the cost $c_k$; they cannot be silently omitted.

If one experimental cycle takes $t_{{\rm cyc},k}$ and has acceptance probability $p_k$, then a stationary approximation may bind wall-clock integration time $T_k$ to effective trials through
\begin{equation}
 e_k\simeq
 \frac{p_k T_k}{t_{{\rm cyc},k}},
\label{eq:shot-rate}
\end{equation}
provided inter-trial correlations are negligible on the relevant scale. Correlations or drift require a covariance model rather than a nominal trial count.

\subsection{Spectral-density information rates}

For a stationary Gaussian frequency-domain measurement, a common one-sided spectral convention yields an information-rate block of the schematic form
\begin{equation}
 \dot F_{ij,k}
 =
 4\,\mathrm{Re}\int_{f\in\mathcal B_k}
 \mathrm df\,
 \partial_i\tilde\mu_k(f)^\dagger
 S_k^{-1}(f)
 \partial_j\tilde\mu_k(f),
\label{eq:psd-fisher}
\end{equation}
where $S_k(f)$ is the declared noise spectral-density matrix and
$\mathcal B_k$ is the retained analysis band. Under stationary independent accumulation,
\begin{equation}
 F_k(T_k)=T_k\dot F_k.
\end{equation}
The exact numerical prefactor depends on Fourier and one-sided/two-sided spectral conventions; it must be frozen with the likelihood implementation. Equation~\eqref{eq:psd-fisher} is therefore a binding template, not a universal detector specification.

Crucially, the nuisance blocks of $\dot F_k$ must be accumulated together with the target block. Replacing $\dot F_k$ by its target-target element before profiling can overstate the information rate whenever calibration drift, spectral tilt, gain, background, phase, or transfer-function parameters are correlated with the target.

\subsection{Coherence and interrogation time}

Coherence is generally not an additive exposure variable. Let $\tau$ be an interrogation or free-evolution time. It can change both the per-shot response and the shot rate:
\begin{equation}
 \bm s_k=\bm s_k(\tau),
 \qquad
 J_k=J_k(\tau),
 \qquad
 e_k(T,\tau)\simeq
 \frac{p_k(\tau)T}{t_{{\rm cyc},k}(\tau)}.
\label{eq:coherence-map}
\end{equation}
Decoherence can suppress the score at long $\tau$, while control sequences can reshape frequency sensitivity. Quantum-sensing analyses explicitly show that Fisher-information scaling can change when finite coherence time becomes relevant \cite{Gefen2017}. Accordingly, RQIR does not credit ``coherence time'' by simply multiplying Fisher information by a larger number. It is a design variable that must enter the response map.

\subsection{A dimensioned atom-interferometer binding example}

A light-pulse atom interferometer provides a useful first physical binding because its leading acceleration phase has the familiar scale
\begin{equation}
 \phi_k = k_{\rm eff} a T_k^2+\phi_0,
\label{eq:ai-phase}
\end{equation}
where $T_k$ is the pulse-separation time and $\phi_0$ denotes a setting-independent phase offset. Long-baseline atom-interferometer analyses use the same $k_{\rm eff}aT^2$ scale factor and explicitly show that preparation time contributes to wall-clock instability even though it contributes no interferometric phase information \cite{Schilling2020}.

Let the single-shot phase variance be $\sigma_{\phi,k}^2$ and the cycle time be
\begin{equation}
 t_{{\rm cyc},k}=2T_k+T_p.
\end{equation}
For wall-clock allocation $x_k$, the number of independent phase estimates is approximately $x_k/t_{{\rm cyc},k}$, so define the phase-information rate
\begin{equation}
 r_k=\frac{1}{t_{{\rm cyc},k}\sigma_{\phi,k}^2}.
\label{eq:ai-rate}
\end{equation}
If the acceleration $a$ is the target and $\phi_0$ is a free shared nuisance, the two-setting profiled Fisher information is
\begin{equation}
 \mathcal I_a(x_1,x_2)
 =
 \frac{W_1W_2}{W_1+W_2}
 \left[k_{\rm eff}(T_1^2-T_2^2)\right]^2,
 \qquad
 W_k=r_kx_k.
\label{eq:ai-profile}
\end{equation}
A single interrogation time has zero information on $a$ when $\phi_0$ is completely free: acceleration and constant phase are then exactly degenerate. Two distinct $T$ values rotate the target direction relative to the offset nuisance and self-calibrate it.

For a fixed wall time $B=x_1+x_2$, Eq.~\eqref{eq:ai-profile} is maximized by
\begin{equation}
 \frac{x_1^\star}{B}
 =
 \frac{\sqrt{r_2}}{\sqrt{r_1}+\sqrt{r_2}},
 \qquad
 \frac{x_2^\star}{B}
 =
 \frac{\sqrt{r_1}}{\sqrt{r_1}+\sqrt{r_2}}.
\label{eq:ai-allocation}
\end{equation}
Using the realistic staging values $T_p=3\,{\rm s}$ and
$\sigma_\phi=10\,{\rm mrad}$ employed in the Hannover 10-m atom-interferometer sensitivity discussion \cite{Schilling2020}, together with
$T_1=0.4\,{\rm s}$ and $T_2=1.2\,{\rm s}$, gives
\begin{equation}
 (x_1^\star/B,x_2^\star/B)
 =
 (0.45619,0.54381).
\end{equation}
Thus the wall-clock optimum is not an equal split. The longer interrogation receives more time despite its slower cycle because the joint fit values both scale-factor leverage and nuisance separation.

For an explicitly illustrative $k_{\rm eff}=1.6\times10^7\,{\rm m}^{-1}$, the asymptotic information rate from Eq.~\eqref{eq:ai-profile} corresponds to
\begin{equation}
 \sigma_a\sqrt{T_{\rm int}}
 \simeq 2.09\times10^{-9}\ {\rm m\,s^{-2}\sqrt{s}}.
\label{eq:ai-numeric}
\end{equation}
This last number is an RQIR staging calculation, not a reported Hannover instrument performance: the published $T_p$ and phase-noise scale are used only to make the resource units realistic, while the chosen two-setting nuisance model and illustrative $k_{\rm eff}$ are ours. Its purpose is to demonstrate an end-to-end conversion from pulse timing and phase noise to wall-clock profiled information.

\subsection{Geometry and control budgets}

Geometry variables---baseline, source separation, mass placement, arm length, trap distance, detector orientation, or spectral band---usually alter the vectors $\bm s_k$ and $J_k$ themselves. Control resources---pulse complexity, field stability, temperature, isolation, calibration cadence, or state-preparation fidelity---can alter both the score geometry and its cost.

We therefore represent a general apparatus design by controls $\bm u$ and scalable allocations $\bm e$:
\begin{equation}
 (\bm u,\bm e)
 \longmapsto
 \left\{
 \bm s_k(\bm u),J_k(\bm u),C_k(\bm u),c_k(\bm u)
 \right\}.
\label{eq:apparatus-map}
\end{equation}
For fixed $\bm u$, the inner allocation problem can inherit the convex structure of Sec.~\ref{sec:allocation}. Optimizing over $\bm u$ need not be convex and may require a discrete or nonlinear outer search. This nested structure keeps a useful theorem layer without pretending that complete apparatus design is globally convex.

\subsection{Resource-closure criterion}

A branch is \emph{resource-closed at target $I_\star$} only relative to a declared physical design domain $\mathcal R_{\rm phys}$, cost model, and nuisance/covariance model. The criterion is
\begin{equation}
 \exists\,(\bm u,\bm e)\in\mathcal R_{\rm phys}
 \quad\text{such that}\quad
 C(\bm u,\bm e)\le B_{\rm avail},
 \qquad
 \Ires(\bm u,\bm e)\ge I_\star.
\label{eq:resource-closure}
\end{equation}
Failure can occur for two qualitatively different reasons:
\begin{enumerate}
 \item \emph{geometry failure}: the best achievable $\kappa_\star$ vanishes or the target remains nuisance-aligned throughout the admissible design family;
 \item \emph{budget failure}: positive information exists, but the minimum feasible cost exceeds the declared available budget or violates coherence/control constraints.
\end{enumerate}
RQIR records these separately because only the second can be repaired by ``more resources'' without changing the measurement concept.

\section{RQIR-RES-001 resource-conversion certificate}
\label{sec:certificate}

The initial Paper-III implementation is frozen by the source-agnostic RQIR-RES-001 certificate. The deterministic core uses the RQIR-II two-band example; randomized tests probe the general resource map. The reference script uses seed 20260909 and 5000 randomized problems for each of homogeneity, monotonicity, and concavity.

\begin{table*}[t]
\caption{Initial RQIR-RES-001 resource-conversion regressions. Randomized tests use fixed seed 20260909.}
\label{tab:resource-certificate}
\begin{ruledtabular}
\begin{tabular}{p{0.18\textwidth}p{0.49\textwidth}p{0.24\textwidth}c}
Test & Invariant or expected result & Reference result & Status\\
\hline
A. Homogeneity & $\mathcal I(\alpha\bm e;0)=\alpha\mathcal I(\bm e;0)$ & max error $2.84\times10^{-13}$ over 5000 trials & PASS\\
B. Monotonicity & $\Delta\bm e\succeq0\Rightarrow\mathcal I(\bm e+\Delta\bm e)\ge\mathcal I(\bm e)$ & minimum gain $3.77\times10^{-2}$ over 5000 trials & PASS\\
C. Concavity & $\mathcal I(t\bm e_1+(1-t)\bm e_2)\ge t\mathcal I(\bm e_1)+(1-t)\mathcal I(\bm e_2)$ & minimum gap $5.50\times10^{-6}$ over 5000 trials & PASS\\
D. Alignment obstruction & common exact nuisance alignment remains zero under resource scaling & max $|\mathcal I|=1.16\times10^{-10}$ through budget $10^6$ & PASS\\
E. Two-band optimum & analytic $(1.4,0.6)$ at $B=2$ & grid optimum $e_2=1.4000$, $\mathcal I=0.3528$ & PASS\\
F. Budget inversion & $B_{\min}=I_\star/\kappa_\star$ & $\sigma_\star=0.2\Rightarrow B_{\min}=141.7233560$, reconstructed $I=25$ & PASS\\
G. Envelope gradient & Eq.~\eqref{eq:marginal} agrees with finite differences & max error $8.95\times10^{-8}$ & PASS\\
H. Fixed-prior separation & fixed $\Lambda$ breaks exposure homogeneity while information remains monotone & $I(10\bm e)-10I(\bm e)=-1.05820$ & PASS\\
I. Atom wall-time binding & two interrogation times self-calibrate a free constant phase offset & $(x_1,x_2)/B=(0.45619,0.54381)$; $\sigma_a\sqrt{s}=2.09\times10^{-9}\,\mathrm{m\,s^{-2}\sqrt{s}}$ for illustrative $k_{\rm eff}$ & PASS\\
\end{tabular}
\end{ruledtabular}
\end{table*}

The randomized tests do not establish a new theorem; Propositions 1 and 2 are analytic. Their purpose is to make implementation regressions visible before an apparatus-specific resource model is attached. In particular, Test H prevents a fixed external nuisance prior from being accidentally multiplied with science exposure, while Test D carries forward the Paper-II rule that exact nuisance alignment is not an exposure shortage.

\subsection{What the certificate still does not test}

RQIR-RES-001 presently assumes that each declared setting has a fixed per-unit local score model and that the scalable resource enters linearly in the Fisher blocks. It does not validate detector dead-time models, colored nonstationary drift, nonlinear likelihood curvature, quantum-measurement attainability, integer scheduling, apparatus engineering limits, or a specific gravity-interface signal amplitude. Those belong to the binding layer for each candidate experiment.

A later apparatus-specific certificate should therefore contain, at minimum, the following frozen objects: the physical parameter/reference point, response Jacobians, covariance or spectral-density model, nuisance-sharing graph, calibration-prior provenance, mapping from wall-clock or shot resources to accepted information, resource costs, feasibility constraints, and an independently reproducible budget inversion.
